\chapter{DF-064 Audit keybinding help and shell hints family across real surfaces}\label{record-105}

\noindent{\footnotesize\color{BookMuted}\path{docs/design/DF-064-audit-keybinding-help-and-shell-hints-family-across-real-surfaces.md}\par}\smallskip

\begin{RecordMeta}
title: DF-064 Audit keybinding help and shell hints family across real surfaces
legend: DF
lane: design
priority: medium
keywords:
  - component-audit
  - dogfood
  - lowerings
  - help

\end{RecordMeta}

\section{Framing}\label{record-105-framing}

Keybinding help is shell infrastructure, not decorative footer copy. The family starts with \texttt{c\allowbreak{}r\allowbreak{}e\allowbreak{}a\allowbreak{}t\allowbreak{}e\allowbreak{}K\allowbreak{}e\allowbreak{}y\allowbreak{}M\allowbreak{}a\allowbreak{}p\allowbreak{}(\allowbreak{})} as structured shortcut truth, then lowers that truth through full grouped references (\texttt{h\allowbreak{}e\allowbreak{}l\allowbreak{}p\allowbreak{}V\allowbreak{}i\allowbreak{}e\allowbreak{}w\allowbreak{}(\allowbreak{})} / \texttt{h\allowbreak{}e\allowbreak{}l\allowbreak{}p\allowbreak{}V\allowbreak{}i\allowbreak{}e\allowbreak{}w\allowbreak{}S\allowbreak{}u\allowbreak{}r\allowbreak{}f\allowbreak{}a\allowbreak{}c\allowbreak{}e\allowbreak{}(\allowbreak{})}), compact shell hints (\texttt{h\allowbreak{}e\allowbreak{}l\allowbreak{}p\allowbreak{}S\allowbreak{}h\allowbreak{}o\allowbreak{}r\allowbreak{}t\allowbreak{}(\allowbreak{})} / \texttt{h\allowbreak{}e\allowbreak{}l\allowbreak{}p\allowbreak{}S\allowbreak{}h\allowbreak{}o\allowbreak{}r\allowbreak{}t\allowbreak{}S\allowbreak{}u\allowbreak{}r\allowbreak{}f\allowbreak{}a\allowbreak{}c\allowbreak{}e\allowbreak{}(\allowbreak{})}), and filtered views (\texttt{h\allowbreak{}e\allowbreak{}l\allowbreak{}p\allowbreak{}F\allowbreak{}o\allowbreak{}r\allowbreak{}(\allowbreak{})} / \texttt{h\allowbreak{}e\allowbreak{}l\allowbreak{}p\allowbreak{}F\allowbreak{}o\allowbreak{}r\allowbreak{}S\allowbreak{}u\allowbreak{}r\allowbreak{}f\allowbreak{}a\allowbreak{}c\allowbreak{}e\allowbreak{}(\allowbreak{})}).

This audit verifies the existing DOGFOOD \texttt{h\allowbreak{}e\allowbreak{}l\allowbreak{}p\allowbreak{}-\allowbreak{}v\allowbreak{}i\allowbreak{}e\allowbreak{}w} story against real story rendering, component-family doctrine, and mode-aware lowerings. It does not redesign keymaps; it turns the current help and hint contract into executable release proof.

\section{Sponsored Users}\label{record-105-sponsored-users}

\begin{itemize}
\tightlist
\item
  TUI app authors building keyboard-owned shells that need shortcut discovery without turning the command palette into a help page.
\item
  Docs readers checking how compact hints and grouped help survive without borders, color, or rich surface chrome.
\item
  Maintainers who need one focused regression tying help story metadata, variants, and component-family guidance to real output.
\end{itemize}

\section{Hills}\label{record-105-hills}

\begin{enumerate}
\def\labelenumi{\arabic{enumi}.}
\tightlist
\item
  A builder can inspect DOGFOOD and understand when to show a compact shell hint versus a grouped shortcut reference.
\item
  A reader can switch both help variants through interactive, static, pipe, and accessible profiles while preserving shortcut grouping and action labels.
\item
  A maintainer can run one cycle test that proves help and shell-hint story metadata and design-system guidance have not drifted away from runtime behavior.
\end{enumerate}

\section{Playback Questions}\label{record-105-playback-questions}

\begin{itemize}
\tightlist
\item
  Does DOGFOOD expose the family through the \texttt{h\allowbreak{}e\allowbreak{}l\allowbreak{}p\allowbreak{}-\allowbreak{}v\allowbreak{}i\allowbreak{}e\allowbreak{}w} story?
\item
  Do variants cover both compact shell hints and grouped shortcut reference?
\item
  Do rich/static previews keep structured surface rendering on the Surface path?
\item
  Do pipe lowerings keep shortcut summaries and grouped help blocks readable without box chrome?
\item
  Do accessible lowerings preserve group names, shortcut keys, action labels, and shell scope in reading order?
\item
  Do component-family docs still separate help discovery from command execution and status/event messaging?
\end{itemize}

\section{Requirements}\label{record-105-requirements}

\begin{itemize}
\tightlist
\item
  Keep the existing \texttt{h\allowbreak{}e\allowbreak{}l\allowbreak{}p\allowbreak{}-\allowbreak{}v\allowbreak{}i\allowbreak{}e\allowbreak{}w} story identity and \texttt{k\allowbreak{}e\allowbreak{}y\allowbreak{}b\allowbreak{}i\allowbreak{}n\allowbreak{}d\allowbreak{}i\allowbreak{}n\allowbreak{}g\allowbreak{}-\allowbreak{}h\allowbreak{}e\allowbreak{}l\allowbreak{}p\allowbreak{}-\allowbreak{}a\allowbreak{}n\allowbreak{}d\allowbreak{}-\allowbreak{}s\allowbreak{}h\allowbreak{}e\allowbreak{}l\allowbreak{}l\allowbreak{}-\allowbreak{}h\allowbreak{}i\allowbreak{}n\allowbreak{}t\allowbreak{}s} coverage family id.
\item
  Render every help story variant in every canonical story profile.
\item
  Treat box drawing as visual-profile chrome only.
\item
  Keep pipe and accessible lowerings shortcut-label-first and semantic-preserving.
\item
  Confirm \texttt{d\allowbreak{}o\allowbreak{}c\allowbreak{}s\allowbreak{}/\allowbreak{}d\allowbreak{}e\allowbreak{}s\allowbreak{}i\allowbreak{}g\allowbreak{}n\allowbreak{}-\allowbreak{}s\allowbreak{}y\allowbreak{}s\allowbreak{}t\allowbreak{}e\allowbreak{}m\allowbreak{}/\allowbreak{}c\allowbreak{}o\allowbreak{}m\allowbreak{}p\allowbreak{}o\allowbreak{}n\allowbreak{}e\allowbreak{}n\allowbreak{}t\allowbreak{}-\allowbreak{}f\allowbreak{}a\allowbreak{}m\allowbreak{}i\allowbreak{}l\allowbreak{}i\allowbreak{}e\allowbreak{}s\allowbreak{}.\allowbreak{}m\allowbreak{}d} still matches the runtime story posture.
\end{itemize}

\section{Acceptance Criteria}\label{record-105-acceptance-criteria}

\begin{itemize}
\tightlist
\item
  \texttt{t\allowbreak{}e\allowbreak{}s\allowbreak{}t\allowbreak{}s\allowbreak{}/\allowbreak{}c\allowbreak{}y\allowbreak{}c\allowbreak{}l\allowbreak{}e\allowbreak{}s\allowbreak{}/\allowbreak{}D\allowbreak{}F\allowbreak{}-\allowbreak{}0\allowbreak{}6\allowbreak{}4\allowbreak{}/\allowbreak{}k\allowbreak{}e\allowbreak{}y\allowbreak{}b\allowbreak{}i\allowbreak{}n\allowbreak{}d\allowbreak{}i\allowbreak{}n\allowbreak{}g\allowbreak{}-\allowbreak{}h\allowbreak{}e\allowbreak{}l\allowbreak{}p\allowbreak{}-\allowbreak{}f\allowbreak{}a\allowbreak{}m\allowbreak{}i\allowbreak{}l\allowbreak{}y\allowbreak{}-\allowbreak{}a\allowbreak{}u\allowbreak{}d\allowbreak{}i\allowbreak{}t\allowbreak{}.\allowbreak{}t\allowbreak{}e\allowbreak{}s\allowbreak{}t\allowbreak{}.\allowbreak{}t\allowbreak{}s} proves the cycle doc carries the modern playback sections.
\item
  The test locates the DOGFOOD \texttt{h\allowbreak{}e\allowbreak{}l\allowbreak{}p\allowbreak{}V\allowbreak{}i\allowbreak{}e\allowbreak{}w\allowbreak{}(\allowbreak{})\allowbreak{} \allowbreak{}/\allowbreak{} \allowbreak{}h\allowbreak{}e\allowbreak{}l\allowbreak{}p\allowbreak{}S\allowbreak{}h\allowbreak{}o\allowbreak{}r\allowbreak{}t\allowbreak{}S\allowbreak{}u\allowbreak{}r\allowbreak{}f\allowbreak{}a\allowbreak{}c\allowbreak{}e\allowbreak{}(\allowbreak{})} story and verifies expected variants.
\item
  Every variant renders non-empty output in every documented profile.
\item
  Rich/static output includes structured visual containment; pipe and accessible output do not depend on box drawing.
\item
  Pipe and accessible output preserve compact shell hints, grouped help title, group names, shortcut keys, and action labels.
\item
  Component-family docs retain \texttt{c\allowbreak{}r\allowbreak{}e\allowbreak{}a\allowbreak{}t\allowbreak{}e\allowbreak{}K\allowbreak{}e\allowbreak{}y\allowbreak{}M\allowbreak{}a\allowbreak{}p\allowbreak{}(\allowbreak{})}, help view, short hint, filtered help, shell-scope, pipe, and accessible guidance.
\end{itemize}

\section{Implementation Outline}\label{record-105-implementation-outline}

\begin{enumerate}
\def\labelenumi{\arabic{enumi}.}
\tightlist
\item
  Add a DF-064 cycle test that renders the \texttt{h\allowbreak{}e\allowbreak{}l\allowbreak{}p\allowbreak{}-\allowbreak{}v\allowbreak{}i\allowbreak{}e\allowbreak{}w} DOGFOOD story through the shared story protocol.
\item
  Read-test the keybinding help and shell hints section in \texttt{d\allowbreak{}o\allowbreak{}c\allowbreak{}s\allowbreak{}/\allowbreak{}d\allowbreak{}e\allowbreak{}s\allowbreak{}i\allowbreak{}g\allowbreak{}n\allowbreak{}-\allowbreak{}s\allowbreak{}y\allowbreak{}s\allowbreak{}t\allowbreak{}e\allowbreak{}m\allowbreak{}/\allowbreak{}c\allowbreak{}o\allowbreak{}m\allowbreak{}p\allowbreak{}o\allowbreak{}n\allowbreak{}e\allowbreak{}n\allowbreak{}t\allowbreak{}-\allowbreak{}f\allowbreak{}a\allowbreak{}m\allowbreak{}i\allowbreak{}l\allowbreak{}i\allowbreak{}e\allowbreak{}s\allowbreak{}.\allowbreak{}m\allowbreak{}d}.
\item
  Move the backlog note into \texttt{d\allowbreak{}o\allowbreak{}c\allowbreak{}s\allowbreak{}/\allowbreak{}d\allowbreak{}e\allowbreak{}s\allowbreak{}i\allowbreak{}g\allowbreak{}n\allowbreak{}/} and update the v6 lane pointer.
\item
  Record any drift as follow-on backlog debt instead of widening this audit.
\end{enumerate}

\section{Drift Check}\label{record-105-drift-check}

No runtime drift was found in this slice. The existing \texttt{h\allowbreak{}e\allowbreak{}l\allowbreak{}p\allowbreak{}-\allowbreak{}v\allowbreak{}i\allowbreak{}e\allowbreak{}w} story already distinguishes compact shell hints from grouped shortcut reference, renders structured surfaces in interactive/static profiles, and lowers to plain shortcut summaries plus grouped help text in pipe/accessibility modes.

No follow-on backlog item is needed for this slice.

\section{Playback}\label{record-105-playback}

\begin{itemize}
\tightlist
\item
  RED: the release lane had only a backlog note; no DF-064 playback test tied keybinding help to every canonical profile.
\item
  GREEN: the new cycle test renders both \texttt{s\allowbreak{}h\allowbreak{}e\allowbreak{}l\allowbreak{}l\allowbreak{}-\allowbreak{}h\allowbreak{}i\allowbreak{}n\allowbreak{}t} and \texttt{g\allowbreak{}r\allowbreak{}o\allowbreak{}u\allowbreak{}p\allowbreak{}e\allowbreak{}d\allowbreak{}-\allowbreak{}r\allowbreak{}e\allowbreak{}f\allowbreak{}e\allowbreak{}r\allowbreak{}e\allowbreak{}n\allowbreak{}c\allowbreak{}e} variants across interactive, static, pipe, and accessible profiles.
\item
  Rich/static previews preserve structured visual containment and shortcut rhythm.
\item
  Pipe previews preserve shortcut summaries and grouped help blocks without box drawing.
\item
  Accessible previews linearize the same shortcut groups and action labels in reading order.
\end{itemize}

\section{Retrospective}\label{record-105-retrospective}

The keybinding help family did not need runtime changes. The useful release work was proving that shortcut discovery is structured shell truth, so future blocks can reuse keymaps without rebuilding help text by hand.

